% Exemple de tableaux





Un  tableau permet de synthétiser des résultats issus de simulations numériques ou des valeurs expérimentales 
ou de fournir des paramètres en un coup d'oeil. 
Il contient donc toujours des informations jugées  essentielles sous la forme de données numériques ou de texte.

\begin{itemize}
\item Il est déconseillé de positionner des tableaux au milieu du texte lorsque cela diminue sa visibilité.
\item Il vaut mieux les insérer comme des figures avec un numéro et une légende.
\item Il ne faut pas oublier de positionner le tableau près de l'endroit où il est évoqué et de le nommer par sa référence
\item Un tableau doit rester sobre.  
 \item On peut ajouter quelques couleurs par alternance, par lignes lorsqu'il y a beaucoup de lignes pour améliorer la lecture
 et minimiser les erreurs de lecture. On peut aussi procéder ainsi pour les colonnes.
\item Préférer un tableau avec des colonnes larges que des colonnes serrées et il faut  centrer le tableau sur la page.
\item le tableau devra être amélioré en ajoutant des couleurs et en augmentant les tailles de fontes 
pour une présentation en mode diapositive
\item s'il y a beaucoup d'information à faire passer, il faut mieux proposer plusieurs tableaux qu'un tableau illisible ou incompréhensible
\end{itemize}




\begin{table}[htb]

\vs{0.5}

\setstretch{1.4}
\small
\centering
\begin{tabular}{p{1.5cm}p{2.5cm}p{2.5cm}p{3.5cm}p{3cm}p{2cm} }
\Xhline{2pt}
Case    &      Fluid    & Velocity     & Reynolds number   &  Flow rate     & Length \\ 
        &               & $U$~[m/s]    & $Re$ in $10^6$  &  $q_m$~[kg/m$^3$] &  $\ell$~[mm] \\ \Xhline{1pt}
1       & Air           & $100$        &   $1.5$         &  $100$            &  $35.33$     \\
2       & Water         &  $\;\;10$        &   $0.1$         &  $\;\;78$             &  $\;\;5.24$  \\
\Xhline{2pt}
\end{tabular}
\caption{Premier exemple  de format de tableau (excellent).}           
\label{tab:exemple1}   
\normalsize
\setstretch{1}
\end{table}


Les tableaux tabs.~\ref{tab:exemple1}, \ref{tab:exemple2} et \ref{tab:exemple3} sont trois exemples qu'on peut rencontrer.

Commentons les.



\begin{enumerate}
\item Le premier tableau tab.~\ref{tab:exemple1} est rencontré dans les publications scientifiques. Il respecte les consignes suivantes:
		\begin{itemize}
		\item Il est sobre avec deux lignes épaisses pour délimiter le tableau dans le plan vertical et une ligne fine
		 qui démarque le titre des colonnes
		 \item Les grandeurs sont clairement décrites avec leur symbôle et leur unités physiques 
		 \item l'interlignage ou espacement vertical est suffisamment important pour améliorer la lisibilité
		 \item les valeurs ou textes sont positionnés à gauche et alignés (on peut aligner sur le . pour les nombres réels)
		 \item le tableau occupe presque la largeur de la page, les colonnes sont dimensionnées précisément.
		 \item le tableau est bien séparé du texte et ressort bien
		 \item on peut remarquer que la hauteur de fonte est légèrement inférieure à celle du texte principale, ce qui peut
		 permettre d'ajouter plus de colonnes et de lignes sans occuper trop de place.
		\end{itemize}
		
		C'est la référence à suivre.
		

		
\item Le second tableau tab.~\ref{tab:exemple1} est aussi acceptable mais souffre de quelques défauts:
		\begin{itemize}
		\item on ne connaît pas le sens des variables
		\item les unités sont répétées inutilement
		\item les colonnes sont trop serrées, elles n'occupent pas l'espace en largeur
		\end{itemize}	
		on peut ajouter qu'il n'y a pas de délimitation horizontale et que l'interlignage est très fort. Mais celà
		n'est pas nécessairement négatif.
		
		On peut alors rajouter le sens des variables dans la légende qu'il faut bien lire pour comprendre le tableau.
		
\begin{table}[htb]
\centering
\setstretch{2}
\begin{tabular}{c|cccc}
Case             & Fluid    & $U$               & $Re$                   &   $q_m$ \\ \hline\hline
1                & Air     & $100$~m/s        &   $1.5 \times 10^6$    &  $100$~kg/m$^3$                       \\
2                & Water   &  $10$~m/s        &   $0.1 \times 10^6$    &   $5$ ~kg/m$^3$      \\
\end{tabular}
\caption{Deuxième exemple de format de tableau (acceptable). La vitesse est notée $U$ et le débit massique $q_m$. }           
\label{tab:exemple2}   
\normalsize
\setstretch{1}
\end{table}
		
\item Le dernier exemple  tab.~\ref{tab:exemple1} est une représentation très basique des données sans aucune réflexion sur le rendu. Ils souffrent de tous les défauts possibles
	   qui sont le contraire des qualités du premier tableau. Il n'est en particulier pas nécessaire de marquer les lignes et
	   les colonnes quand c'est évident. Et il n'y ni unité, ni explication sur le sens des variables. Enfin les lignes et les colonnes sont très resserrées.
	   On peut aussi remarquer que le tableau n'est pas centré sur la page.
	   
	   Cette représentation est à proscrire!
	   
	   
\begin{table}[htb]
\begin{tabular}{|c|c|c|c|c|}
\hline
Case             & Fluid    & $U$          & $Re$                   &   $q_m$ \\ \hline
1                & Air     & $100$        &   $1.5 \times 10^6$    &  $100$    \\ \hline
2                & Water   &  $10$        &   $0.1 \times 10^6$    &   $5$      \\ \hline
\end{tabular}
\caption{Troisième exemple de format de tableau (le plus mauvais) }           
\label{tab:exemple3}   
\normalsize
\setstretch{1}
\end{table}
	   
\end{enumerate}


On peut aussi:

\begin{itemize}
\item ajouter de la couleur soit par ligne, soit par colonne, soit une ligne sur deux et une colonne sur deux pour encore améliorer la visibilité.
\item choisir une couleur unique avec une teinte foncée et claire pour éviter l'effet "arc-en-ciel"
\item encore jouer sur les épaisseurs des traits.
\item modifier les types de fontes par rapport à celle du texte principal pour insister sur le contenu des colonnes
\end{itemize}





